% !TeX root = ../教材题库_全部条目_转写工作稿.tex
% Source: XJ-B1-C06-RV-P265-Y0217
% 第6章 / 章末复习 复习题六 / 教材 p.265 / PDF 第272页
% OCR未校对：公式与文字均可能有误；未完成数学TeX转写；禁止直接发给学生
% 原题及相邻原解截图（核对用）：../题目原图/XJ-B1-C06-RV-P265-Y0217-01.webp
% 本文件只在转写完成后作为题干使用；OCR 条目尚不是可靠题干。
\begin{bankitem}{复习题六 第6题}{第6章 / 章末复习 复习题六 / 教材 p.265 / PDF 第272页}{XJ-B1-C06-RV-P265-Y0217}{OCR未校对：公式与文字均可能有误；未完成数学TeX转写；禁止直接发给学生}
\begin{ocrdraft}
[教材p.265 / PDF第272页]
\par
6．一条产品生产线平均每天的产量为3700件，标准差为50件．如果某一
\par
天的产量低于或高于平均产量，并落在\ensuremath{\pm}2个标准差的范围之外，就认为该生产
\par
线失去控制．下表是该生产线一周各天的产量，该生产线哪几天失去了控制？
\par
时间
\par
星期一
\par
星期二
\par
星期三
\par
星期四
\par
星期五
\par
星期
\par
星期日
\par
产量/件
\par
3850
\par
3670
\par
3690
\par
3720
\par
3610
\par
3590
\par
3700
\end{ocrdraft}
\end{bankitem}
